\section{Notation}

\subsection{Tensor networks}

We here use the tensor network notation \cite{goessmann_tensor_2026}.
The modes of each tensors are assigned by variables, and denoted in brackets as
\begin{align*}
    \hypercorewith \, .
\end{align*}
Quantum gates are along this line unitary tensors with incoming and outgoing variables.
The Pauli-X for example is denoted by
\begin{align*}
    \paulixat{\cavariable{\insymbol},\cavariable{\outsymbol}} \coloneqq
    \coloredmatrixof{0 & 1 \\ 1 & 0}{\cavariable{\insymbol},\cavariable{\outsymbol}} \, .
\end{align*}
Contractions are denoted by angle brackets with the variables left open specified in further brackets.
Circuits are contractions of the unitaries, leaving only the last output variables open.
Computational basis measurements of circuits are coordinatewise squared absolut of the contracted circuit.

\subsection{\ComputationActivationNetworks{}}

\ComputationActivationNetworks{} (\CompActNets{}) are tensor networks representing logical, probabilistic and neural models \cite{goessmann_tensor_2026}.
\red{We here introduce them for Boolean statistics.}

To any Boolean function $\formula:[2]^{\catorder}\rightarrow[2]$ we define its basis encoding
\begin{align*}
    \bencodingofat{\formula}{\headvariable,\shortcatvariables}
    = \sum_{\shortcatindices\in[2]^{\catorder}} \onehotmapofat{\formula(\shortcatindices)}{\headvariable} \otimes \onehotmapofat{\shortcatindices}{\shortcatvariables}\, ,
\end{align*}
where we denote the basis vectors as $\fbasisat{\catvariable}=\coloredmatrixof{1 \\ 0}{\catvariable}$ and $\tbasisat{\catvariable}=\coloredmatrixof{0 \\ 1}{\catvariable}$.

\begin{definition}[\ComputationActivationNetwork{} (\CompActNets{})]
    \label{def:compActNets}
    Let there be a function $\sstat : [2]^{\catorder} \allowbreak \rightarrow [2]^{\seldim}$ with basis encoding $\bencodingofat{\sstat}{\headvariables,\shortcatvariables}$.
    Let there further be a hypergraph $\graph=(\nodes,\edges)$ with nodes $\nodes$ containing $[\seldim]$.
    We define the by $\sstat$ computable and by $\graph$ activated family of distributions by
    \begin{align*}
        \realizabledistsof{\sstat,\graph}
        = \left\{ \normalizationof{\bencodingofat{\sstat}{\headvariables,\shortcatvariables},\contractionof{\acttensor}{\headvariables} %\} \cup \{\hypercoreofat{\edge}{\sstatcatof{\edge}} \wcols \edgein \}
        }{\shortcatvariables}
              \wcols \acttensorat{\headvariableof{\nodes}} \in \tnsetof{\graph} \right\} \, .
    \end{align*}
    We refer to any member $\probat{\shortcatvariables}\in\realizabledistsof{\sstat,\graph}$ as a \emph{\ComputationActivationNetwork{}} (or shorter as a \emph{\CompActNet{}}).
    We call $\bencsstatwith$ (and any decomposition of it) the \emph{computation network} and $\acttensorat{\headvariableof{\nodes}}$ the \emph{activation network}.
\end{definition}

\subsection{Q-samples}

In general, we define q-samples to be quantum states, which measured in the computational basis reproduce a given probability distribution.

\begin{definition}[Q-sample \cite{aharonov_adiabatic_2003,low_quantum_2014}]
    Given a probability distribution $\probtensor:\atomstates\rightarrow\rr$ (i.e. $\contraction{\probtensor}=1$ and $\zeros \prec \probtensor$) its q-sample is
    \begin{align*}
        \qstateofat{\probtensor}{\shortcatvariables}
        = \sum_{\shortcatindicesin} \sqrt{\probat{\indexedshortcatvariables}} \cdot \onehotmapofat{\shortcatindices}{\shortcatvariables} \, .
    \end{align*}
\end{definition}

Q-samples are more restrictive than arbitrary states having a distribution as a measurement distribution, since they demand real and positive amplitudes.

\begin{example}[Q-sample of the uniform distribution]
    \label{exa:qSampleUniform}
    The q-sample of the uniform distribution can be prepared by Hadamard gates acting on the ground state.
    This is an example of a Walsh-Hadamard transform, which can be performed by unary Hadamard gates.
\end{example}

Generalizing \exaref{exa:qSampleUniform}, Q-samples to independent distributions of boolean variables can be prepared by unary rotations along the y-axis.
